% grav_universal: the apple and the Moon obey one law (Ch 16, fig:gravuniv)
% Numbers: r = 3.84e8 m = 60.3 R_E, T = 27.32 d = 2.36e6 s,
%          a_Moon = 4 pi^2 r / T^2 = 2.72e-3 m/s^2 = g/3600.
\begin{tikzpicture}[x=1cm,y=1cm]
  \def\D{10.6}  % Earth-Moon distance on the page (not to scale)
  % the Moon's orbit, a short arc of it
  \draw[orbitpath] (-8:\D) arc[start angle=-8,end angle=9,radius=\D];
  % Earth
  \draw[planet] (0,0) circle[radius=0.85];
  \node at (0,0) {Earth};
  % the apple and its acceleration
  \draw[moon] (0,1.55) circle[radius=0.11];
  \draw[accel] (0,1.36) -- (0,0.93);
  \node[cPrim,anchor=west] at (0.18,1.2) {$a=g$};
  \node[anchor=east] at (-0.2,1.6) {apple};
  % Moon, its velocity and its acceleration
  \draw[moon] (\D,0) circle[radius=0.26];
  \node[anchor=west] at (\D+0.32,0) {Moon};
  \draw[velocity] (\D,0.3) -- (\D,1.25) node[anchor=west,xshift=1pt] {$v$};
  \draw[accel] (\D-0.3,0) -- (\D-1.25,0);
  \node[cPrim,anchor=north] at (\D-0.8,-0.08) {$a_{\mathrm{Moon}}$};
  % the distance
  \draw[dimtick] (0,-0.95) -- (0,-1.55);
  \draw[dimtick] (\D,-0.32) -- (\D,-1.55);
  \draw[dim] (0,-1.42) -- node[below=1pt,cSteel] {$r=60\,R_E$ \ (not to scale)} (\D,-1.42);
  % the two measurements
  \node[anchor=north west,align=left] at (-0.95,3.15)
    {\textbf{From dropping things}\\[1pt] $g=\SI{9.8}{m/s^2}$};
  \node[anchor=north west,align=left] at (6.35,3.15)
    {\textbf{From the Moon's period}\\[1pt]
     $a_{\mathrm{Moon}}=\dfrac{4\pi^2 r}{T^2}=\SI{0.0027}{m/s^2}$};
  % the comparison
  \node at (\D/2,-2.45)
    {$\dfrac{a_{\mathrm{Moon}}}{g}=\dfrac{0.0027}{9.8}\approx\dfrac{1}{3600}=\dfrac{1}{60^2}$,
     \ as an inverse-square law requires};
\end{tikzpicture}
